% Generated draft: requires mathematical and source review.
\documentclass[11pt]{article}
\usepackage[margin=25mm]{geometry}
\usepackage{fontspec}
\setmainfont{Latin Modern Roman}
\usepackage{amsmath,amssymb,mathtools,unicode-math}
\setmathfont{Latin Modern Math}
\usepackage{enumitem,needspace,xurl,hyperref}
\setlength{\parindent}{0pt}
\setlength{\parskip}{5pt}
\setlength{\emergencystretch}{3em}
\newcommand{\N}{\mathbb{N}}
\newcommand{\Z}{\mathbb{Z}}
\newcommand{\Q}{\mathbb{Q}}
\newcommand{\R}{\mathbb{R}}
\newcommand{\C}{\mathbb{C}}
\newcommand{\F}{\mathbb{F}}
\newcommand{\A}{\mathbb{A}}
\newcommand{\PP}{\mathbb{P}}
\newcommand{\E}{\mathbb{E}}
\newcommand{\eps}{\varepsilon}
\newcommand{\dd}{\,\mathrm{d}}
\begin{document}
\section*{UnsolvedMath ID: 30004381}
\textbf{Source problem code:} catalogue entry
\subsection{Definitions and mathematical statement}
Let \(d\) be a positive integer and let \(D_d\subset\R^d\) be a bounded convex domain. The supplied record specifies
\[
\overline C^2(D_d)=\left\{f\in C^2(D_d):\ \|f\|_\infty\leq 1,\ \operatorname{Lip}(f)\leq d^{-1/2},\ \operatorname{Lip}(D_\theta f)\leq d^{-1}\right\}.
\]
Here \(C^2(D_d)\) denotes twice continuously differentiable functions, \(\|f\|_\infty=\sup_{x\in D_d}|f(x)|\), and \(D_\theta f\) denotes differentiation in direction \(\theta\). The notation \(\operatorname{Lip}\) denotes a Lipschitz seminorm; the record does not explicitly specify its underlying norm or metric. It also does not specify the range or quantifier for \(\theta\). In particular, requiring the last inequality for every Euclidean unit direction would be an additional interpretation requiring confirmation.

The proposed conjecture is that there exists an absolute constant \(C>0\), independent of dimension and domain, such that
\[
 e_n\bigl(\overline C^2(D_d),\operatorname{INT}\bigr)\geq C n^{-2/d}
\]
for every positive integer \(d\), every admissible \(D_d\), and every positive integer \(n\). Positive integers are the implicit index ranges used here; the supplied statement does not spell them out. The quantity \(e_n\) is described only as the \(n\)th minimal integration error.

This remains an incompletely specified conjecture: the supplied record does not define \(\operatorname{INT}\), its measure or normalization, the information available to algorithms, the meaning of the information budget \(n\), or the error criterion. It does not state any volume or geometric normalization of \(D_d\), nor any boundary extension requirement on \(f\). These choices, together with the directional quantifier, must be recovered from the source before a precise mathematical formulation can be certified; none is imposed here.
\subsection{Short English statement}
The supplied conjecture asks for a lower bound \(C n^{-2/d}\) on the minimal integration error for a class of twice continuously differentiable functions with dimension-dependent Lipschitz bounds, using one constant for all dimensions and admissible convex domains. The directional condition, domain normalization, and integration model require clarification. This is a statement-only draft requiring human mathematical review.
\subsection{Sources}
\begin{itemize}
\item Local UnsolvedMath registry, \texttt{data/unsolvedmath/problems.json}, permanent ID \(30004381\), source problem code \texttt{catalogue entry}, title: Dimension-Uniform Lower Bounds for Twice-Smooth Integration. This attribution refers to the supplied complete registry record.
\item Supplied report citation: \emph{New Perspectives and Computational Challenges in High Dimensions} (2021), Oberwolfach Reports, pp.\ 377--416, DOI \texttt{10.4171/owr/2020/6}. Supplied source and literature-review link: \url{https://doi.org/10.4171/owr/2020/6}. The record describes these reports as documenting mathematical workshops and their open-problem discussions. No individual proposer is supplied; the proposed year is recorded as 2020.
\item Literature assessment checked at \texttt{2026-08-22}: supplied dated curation, not independently verified here. It labels the problem open and describes the dimension- and domain-uniform lower bound as a stated conjecture, sharp for the cube with grid information; it reports that no general proof or counterexample was found. The supplied background explicitly requests independent verification of time-sensitive status and cited claims. Its listed checked source is the report linked above.
\item The registry labels the statement \texttt{corrected\_verified} and reports that the official report was checked and the function class and dimension-uniform conjecture restored. These are supplied verification claims, not verification performed for this draft.
\end{itemize}
\end{document}
